\section{Introduction}

\begin{figure}[!t]
    \centering
    \includegraphics[width=\linewidth]{figures/intro.pdf}
    \caption{Illustration of our motivation. \textsc{RetroGen} conceptualizes abundant expert artifacts as compressed processes, enabling scalable and grounded process supervision for open-ended agents without costly human or teacher model annotation.}
    \label{fig:intro}
    \vspace{-4mm}
\end{figure}

Large language models (LLMs) are increasingly being extended into autonomous agents capable of planning, tool use, and multi-step interaction~\citep{yao2022react, shinn2023reflexion, wang2024survey, qin2024toolllm}. A key factor driving this progress is trajectory-level supervision: training models not only on final answers, but also on intermediate reasoning, actions, observations, and revisions~\citep{zelikman2022star, lightman2024let, shao2024deepseekmath}. 
In verifiable domains such as mathematics, coding, and closed-ended question answering, generating trajectory data at scale is more feasible because the data can be verified using clear signals.
% sampled at scale and optimized using objective signals, including final-answer correctness, execution feedback, and test outcomes.

However, this reliance on objective verification creates a fundamental bottleneck for open-ended, evidence-grounded tasks, such as synthesizing literature reviews~\citep{openscholar,mcdonald2022detect}, composing financial analyst reports grounded in filings~\citep{loukas2021edgar}, or drafting legal judgments~\citep{guha2023legalbench,dai2025laiw}. Unlike closed-ended problems, these tasks lack a singular ground truth that can validate the generated trajectory. Even when a final output appears coherent, it remains difficult to automatically determine whether the selected evidence is sufficient, whether claims are faithfully grounded, and whether the reasoning path follows domain-specific standards of rigor. As a result, scalable process supervision for evidence-grounded long-form generation remains an important open challenge.

Existing approaches typically obtain trajectory supervision by distilling from stronger teacher models or relying on human experts~\citep{liu2023webglm,qin2023webcpm,lightman2024let}. Yet, both sources are inherently unscalable: teacher distillation requires repeated access to costly proprietary models, while human annotation is prohibitively expensive for evidence-intensive workflows. Moreover, relying on forward trajectory generation in open-ended domains often leads to wandering search behaviors, as models lack the holistic foresight required to synthesize complex information.

This motivates a more self-improving, target-driven alternative. As illustrated in Figure~\ref{fig:intro}, we observe that while high-quality open-ended trajectories are exceptionally scarce, high-quality final artifacts produced by domain experts are remarkably abundant. A published literature review or a binding legal judgment is not merely a static output; rather, it is a lossy compression of a latent, multi-step evidence-seeking process. A well-crafted related work section implicitly encodes how an author searched, filtered, compared, and synthesized prior studies into a coherent narrative. Crucially, because these artifacts are rigorously vetted by domain experts (e.g., peer reviewers), they provide strong quality prior for reverse-engineering the processes that produced them.

This perspective transforms expert-curated artifacts into verifiable targets for scalable process supervision. Instead of demanding an external teacher to demonstrate a full forward trajectory, the agent is tasked with reconstructing a plausible evidence-seeking trajectory that arrives at the known expert artifact.

In this paper, we introduce \textsc{RetroGen}, a self-improving framework that operationalizes retrospective process supervision. Given an expert artifact, \textsc{RetroGen} first establishes an evidence environment by retrieving and organizing relevant sources. The agent then reconstructs a trace detailing how the artifact was derived, encompassing evidence selection, claim grounding, and intermediate synthesis. Rather than treating all generated traces as valid supervision, \textsc{RetroGen} employs a rubric-guided verification step that quantitatively scores candidate trajectories across multiple dimensions. By applying a weighted thresholding mechanism, the agent selectively filters reconstructions based on their comprehensive performance in artifact fidelity, evidence faithfulness, and procedural plausibility. This establishes an iterative artifact-to-trajectory-to-agent training loop.

Our contributions are summarized as follows:
\begin{itemize}[leftmargin=*, topsep=1pt, partopsep=0pt, itemsep=0pt, parsep=3pt] 
    \item We formulate expert-curated artifacts as a scalable source of retrospective process supervision, reframing static outputs as verifiable targets for trajectory reconstruction.
    \item We propose \textsc{RetroGen}, an iterative framework that autonomously generates, verifies, and trains on reconstructed trajectories without relying on stronger teacher models.
    % \item The proposed \textsc{RetroGen} significantly improves evidence-grounded long text generation tasks on both static and agentic benchmarks, while preserving general capabilities.
\end{itemize}

The proposed \textsc{RetroGen} significantly improves evidence-grounded long text generation on both static and agentic benchmarks, while preserving general capabilities. A verification ablation further shows that combining scoring dimensions outperforms any single filter.